% called by main.tex
%
\chapter{Estado de la cuestión}
\label{ch::estado}

Este capítulo sitúa el trabajo en la literatura. Se organiza en cinco áreas que confluyen
en el problema —mercados de predicción, aprendizaje profundo sobre libros de órdenes,
microestructura y ejecución, validación temporal en finanzas, y cambio de
distribución— y termina señalando el hueco concreto que este trabajo pretende cubrir.

\section{Trabajos previos y estudios relevantes}

Esta sección recorre la literatura relevante para el problema, organizada en cinco bloques
temáticos. En cada bloque se resume qué proponen los trabajos de referencia y qué se toma
de ellos para este trabajo. La predicción del movimiento de precios a partir de la
microestructura es un campo activo, con propuestas asentadas tanto para mercados
tradicionales como para criptoactivos; el interés aquí es seleccionar de esa literatura los
elementos aplicables a un mercado de predicción sobre Bitcoin.

\subsection{Mercados de predicción}

Los mercados de predicción se han estudiado, sobre todo, como mecanismos de agregación de
información dispersa. Wolfers y Zitzewitz~\cite{wolfers2004prediction} muestran que pueden
producir previsiones notablemente precisas, si bien su interpretación depende del diseño
del contrato, de la liquidez y de los incentivos de los participantes. La mayor parte de
esta literatura analiza la calibración del precio frente a la resolución final del evento.
El presente trabajo se sitúa un escalón por debajo, en el terreno de la microestructura:
la dinámica intradía del precio del contrato a escala de segundos, donde aparecen
\emph{spreads}, profundidad limitada y desajustes transitorios respecto a las referencias
externas.

\subsection{Aprendizaje profundo sobre libros de órdenes}

La predicción del movimiento de precios a partir del libro de órdenes es un problema con
una literatura ya madura. El trabajo de referencia es DeepLOB~\cite{zhang2019deeplob}, que
combina convoluciones para capturar la estructura espacial del libro con capas recurrentes
para la dependencia temporal, y se valida sobre el conjunto de referencia
FI-2010~\cite{ntakaris2018}. En la misma línea, Tsantekidis et al.~\cite{tsantekidis2017}
emplean redes convolucionales y recurrentes para detectar cambios de precio en el libro.
Sirignano y Cont~\cite{sirignano2019universal} aportan evidencia de la existencia de
características universales en la formación de precios que el aprendizaje profundo es capaz
de capturar, y Jha et al.~\cite{jha2020digital} trasladan estas ideas al terreno de los
criptoactivos. De esta literatura se toma una decisión de diseño importante: tratar el
libro de órdenes como un objeto con estructura espacio-temporal, y no como una tabla
plana. Ahora bien, casi todos estos trabajos se evalúan con métricas de clasificación y
rara vez bajo un modelo de ejecución realista, que es justamente donde este trabajo pone
el acento.

\subsection{Microestructura, costes de ejecución y latencia}

La distancia entre predecir bien y obtener beneficio es uno de los temas clásicos de la
microestructura de mercado. Kearns y Nevmyvaka~\cite{kearns2013} revisan el papel del
aprendizaje automático en microestructura y negociación de alta frecuencia, y subrayan que
la ejecución, el coste y la latencia condicionan de forma decisiva la viabilidad de una
señal. Este trabajo cuantifica ese fenómeno en el contexto concreto de Polymarket: se mide
la latencia real del ciclo de orden y se demuestra que una señal direccional fuerte se
vuelve no rentable bajo una latencia de entrada modesta.

\subsection{Aprendizaje automático en finanzas y validación temporal}

El aprendizaje automático aplicado a la predicción financiera es especialmente propenso a
resultados ilusorios cuando la evaluación no respeta la estructura temporal de los datos.
Gu, Kelly y Xiu~\cite{gu2020empirical} ofrecen un marco riguroso de comparación de modelos
con validación fuera de muestra. López de Prado~\cite{lopezdeprado2018} documenta los
mecanismos por los que se cuela información del futuro y por los que se sobreajustan los
\emph{backtests}, y propone protocolos de validación que respetan la causalidad temporal.
Bailey y López de Prado~\cite{bailey2014deflated} advierten de un riesgo adicional: si se
prueban muchos modelos y se elige el mejor, las métricas de rendimiento salen infladas casi
por construcción. Estas referencias motivan tres decisiones metodológicas de este trabajo:
la partición estrictamente temporal, el protocolo fuera de pliegue (\emph{out-of-fold}) y
la congelación del candidato antes de observar el conjunto de validación fuera de muestra.

\subsection{Cambio de distribución y adaptación de dominio}

Los mercados financieros no son estacionarios, por lo que la distribución de los datos
cambia entre el periodo de entrenamiento y el de despliegue. Para amortiguar ese efecto
existen técnicas de adaptación de dominio, como la alineación de correlaciones
CORAL~\cite{sun2016coral}. En el Capítulo~\ref{ch::resultados} se mostrará un resultado
contraintuitivo: con un horizonte de datos reducido, estas técnicas no superan a una
estrategia mucho más simple, consistente en restringir la operativa al régimen en el que
el modelo fue entrenado. Por último, en cuanto a arquitecturas profundas para series
temporales multivariantes, trabajos como el \emph{Temporal Fusion
Transformer}~\cite{lim2021tft} marcan una línea de evolución natural para cuando se
disponga de muchos más datos.

\section{Limitaciones en la literatura}

Del repaso anterior se desprenden tres limitaciones recurrentes que enmarcan la aportación
de este trabajo. En primer lugar, la evaluación predominante es de \emph{clasificación}
(precisión, AUC), y rara vez se traduce a un resultado económico bajo un modelo de
ejecución realista; como se verá, esa traducción cambia las conclusiones por completo. En
segundo lugar, gran parte de los resultados positivos se obtienen sin una latencia de
entrada explícita, cuando la latencia es precisamente lo que destruye muchas señales de
corto plazo. En tercer lugar, los entornos de criptoactivos jóvenes ofrecen pocos datos,
lo que agrava el riesgo de sobreajuste y de resultados ilusorios si la validación no es
estrictamente temporal.

El hueco que cubre este trabajo es, por tanto, metodológico: evaluar la predicción de
corto plazo \emph{bajo coste y latencia medida}, con validación estrictamente temporal y
congelación previa del candidato, y reportar cada resultado en ticks netos acompañado de
su soporte y su incertidumbre.
